%%%%%%%%%%%%%%%%%%%%%%%%%%%%%%%%%%%%%%%%%%%%%%%%%%%%%%%%%%%%%%%%%%%%%%%%%%%%%%%%

\pagestyle{empty}
\begin{abstract}
  Este trabajo lleva el control visual de coches de tipo Scalextric a una arquitectura
  distribuida. Los trabajos previos resolvían el problema con una aplicación monolítica en la
  que todas las cámaras colgaban por USB de la misma máquina, lo que acota el tamaño del
  circuito que es posible cubrir. Aquí el sistema se descompone en nodos ROS~2 que cooperan a
  través de una red local: un nodo cámara por cada cámara, que detecta la posición de cada
  coche por segmentación de color; un nodo de control por carril, que reconstruye la
  trayectoria durante una vuelta de calibración y decide la tensión que hay que aplicar; y un
  nodo de potencia, que traduce esas consignas en señales PWM sobre el Arduino que alimenta la
  pista. Los nodos pueden repartirse por la red en cualquier combinación, y todo el despliegue
  es reproducible mediante contenedores Docker.
  %, incluido el sistema entero en un solo equipo.

  La calibración es automática y consume una única vuelta, empieza y termina al cruzar la
  meta, sin intervención humana. De ella se obtiene la trayectoria base, sobre la que se apoya
  el resto del sistema. El algoritmo va creando sobre ella zonas de derrape, a las que castiga
  bajando la tensión, mientras que las demás la van subiendo. Cada cámara aporta la parte del
  circuito que ve, con o sin solapamiento.

  Al repartir el lazo de control por la red, la latencia pasa a ser un factor de diseño, por lo
  que los mensajes que lo recorren llevan marcas de tiempo con las que medir cuánto tarda cada
  etapa. Esos mensajes, junto con los de telemetría, los captura un nodo mientras circulan por
  la red para su posterior análisis. El sistema cuenta además con un panel web desde el que
  seguir la carrera en directo.

  El sistema se ha probado en condiciones reales sobre circuitos de distinto tamaño, desde uno
  pequeño que cubre una sola cámara hasta circuitos grandes que necesitan varias cámaras
  repartidas entre varias máquinas, y con diferentes coches. Se comparó el control automático
  con el de pilotos humanos y se estudió la latencia de la red, que queda descartada como
  factor limitante frente a otros elementos del montaje que sí dificultan el control.

  \vspace*{25pt}
  \begin{segundoresumo}
    This project brings the visual speed control of Scalextric cars to a distributed
    architecture. Previous work solved the problem with a monolithic application in which every
    camera hung off the same machine over USB, which bounds the size of the track that can be
    covered. Here the system is broken down into ROS~2 nodes that cooperate over a local
    network: one camera node per physical camera, which detects the position of each car by
    colour segmentation; one control node per lane, which reconstructs the trajectory during a
    calibration lap and decides the voltage to apply; and a power node, which turns those
    setpoints into PWM signals on the Arduino that powers the track. The nodes can be spread
    across the network in any combination, and the whole deployment is reproducible through
    Docker containers.
    %, including running the whole system on a single  machine.

    Calibration is automatic and takes a single lap, starting and ending when the car crosses
    the finish line, with no human intervention. That lap yields the base trajectory, on which
    the rest of the system relies. Over it the algorithm builds skid zones, which it penalises
    by lowering the voltage, while raising it in all the others. Each camera contributes the
    part of the track it sees, with or without overlap.

    Once the control loop is spread over the network, latency becomes a design factor, so the
    messages that travel through it carry timestamps that make it possible to measure how long
    each stage takes. Those messages, together with the telemetry ones, are captured by a node
    as they travel across the network, for later analysis. The system also provides a web
    dashboard from which the race can be followed live.

    The system has been tested under real conditions on tracks of different sizes, from a small
    one covered by a single camera to large tracks that need several cameras spread across
    several machines, and with different cars. The automatic control was compared against human
    drivers, and network latency was studied and ruled out as a limiting factor, unlike other
    elements of the physical setup that do make control harder.
  \end{segundoresumo}
\vspace*{25pt}
\begin{multicols}{2}
\begin{description}
\item [\palabraschaveprincipal:] \mbox{} \\[-20pt]
  \begin{itemize}
    \item Visión artificial
    \item Sistemas distribuidos
    \item ROS~2
    \item Control adaptativo
    \item Seguimiento multicámara
    \item Detección de derrapes
    \item Latencia de red
    \item Scalextric
  \end{itemize}
\end{description}
\begin{description}
\item [\palabraschavesecundaria:] \mbox{} \\[-20pt]
  \begin{itemize}
    \item Artificial vision
    \item Distributed systems
    \item ROS~2
    \item Adaptive control
    \item Multi-camera tracking
    \item Skid detection
    \item Network latency
    \item Scalextric
  \end{itemize}
\end{description}
\end{multicols}

\end{abstract}
\pagestyle{fancy}

%%%%%%%%%%%%%%%%%%%%%%%%%%%%%%%%%%%%%%%%%%%%%%%%%%%%%%%%%%%%%%%%%%%%%%%%%%%%%%%%
